\documentclass[a4paper,13pt,3p,twoside]{report}
\usepackage{scrextend}
\changefontsizes{13pt}
\usepackage[utf8]{vietnam}
\usepackage[top=2cm, bottom=2cm, left=3.5cm, right=2.5cm]{geometry}
\usepackage{xurl}
\usepackage{appendix}
\usepackage{babel}
\usepackage{xcolor}
\usepackage{outlines}
\usepackage{graphicx} % Cho phép chèn hỉnh ảnh
\usepackage{fancybox} % Tạo khung box
\usepackage{indentfirst} % Thụt đầu dòng ở dòng đầu tiên trong đoạn
\usepackage{amsthm} % Cho phép thêm các môi trường định nghĩa
\usepackage{latexsym} % Các kí hiệu toán học
\usepackage{amsmath} % Hỗ trợ một số biểu thức toán học
\usepackage{amssymb} % Bổ sung thêm kí hiệu về toán học
\usepackage{amsbsy} % Hỗ trợ các kí hiệu in đậm
\usepackage{times} % Chọn font Time New Romans
\usepackage{array} % Tạo bảng array
\usepackage{enumitem} % Cho phép thay đổi kí hiệu của list
\usepackage{subfiles} % Chèn các file nhỏ, giúp chia các chapter ra nhiều file hơn
\usepackage{titlesec} % Giúp chỉnh sửa các tiêu đề, đề mục như chương, phần,..
\usepackage{titletoc}
\usepackage{chngcntr} % Dùng để thiết lập lại cách đánh số caption,..
\usepackage{pdflscape} % Đưa các bảng có kích thước đặt theo chiều ngang giấy
\usepackage{afterpage}
\usepackage[ruled,vlined]{algorithm2e}  % Hỗ trợ viết các giải thuật
\usepackage{capt-of} % Cho phép sử dụng caption lớn đối với landscape page
\usepackage{multirow} % Merge cells
\usepackage{fancyhdr} % Cho phép tùy biến header và footer
% \usepackage[natbib,backend=biber,style=ieee]{biblatex} % Giúp chèn tài liệu tham khảo
\usepackage{rotating}
\usepackage[font=small,labelfont=bf]{caption}

\usepackage{listings}
\definecolor{vscBlue}{HTML}{0000FF}      % Xanh dương thuần khiết (Keywords)
\definecolor{vscTeal}{HTML}{00709F}      % Xanh lơ đậm (Classes/Types) 
\definecolor{vscRed}{HTML}{AF0000}       % Đỏ đậm (Strings) 
\definecolor{vscGreen}{HTML}{008000}     % Xanh lá cây (Comments)
\definecolor{vscAnnotation}{HTML}{797980} % Xám trầm (Annotation)
\definecolor{vscNumbers}{HTML}{000000}   % Xám số dòng
\definecolor{vscBg}{HTML}{FFFFFF}        % Nền trắng sạch

\lstset{
    language=Java,
    backgroundcolor=\color{vscBg},
    basicstyle=\bfseries\ttfamily\footnotesize\color{black}, 
    columns=flexible,
    keepspaces=true,
    breakatwhitespace=false,
    breaklines=true,
    captionpos=b,    
    showstringspaces=false,
    frame=single,
    classoffset=0,
    numbers=left,
    numbersep=8pt,
    numberstyle=\tiny\color{vscNumbers},
    commentstyle=\color{vscGreen},
    stringstyle=\color{vscRed},
    keywordstyle=\color{vscBlue}, 
    emph={Application, Scene, Button, StackPane, Stage, JavaFXHello, String, ActionEvent, System, out, println, launch, setText, setOnAction, add, setScene, setTitle, show, getChildren, root, scene, primaryStage, btnHello, evt},
    emphstyle=\color{vscTeal},
    literate={->}{{{\color{vscBlue}->}}}2 
             {@Override}{{{\color{vscAnnotation}@Override}}}9
             {.}{{{\color{black}.}}}1, 
    rulecolor=\color{black},
    showspaces=false,
    showstringspaces=false,
    showtabs=false,
    tabsize=2,
    xleftmargin=0pt,
    xrightmargin=0pt,
    linewidth=\linewidth
}

\usepackage{float}
\usepackage{subcaption}


\usepackage{tcolorbox}
\usepackage{xparse}
\RenewDocumentCommand{\texttt}{v}{%
  \tcbox[
    boxrule=0.4pt,
    arc=1pt,
    colback=gray!15,
    colframe=gray!30,
    coltext=black,
    left=2pt, right=2pt,
    top=1pt, bottom=1pt,
    boxsep=0pt,
    fontupper=\bfseries\ttfamily\footnotesize,
    valign=center,
    nobeforeafter,
    box align=base
  ]{#1}%
}


\usepackage[nonumberlist, nopostdot, nogroupskip, acronym]{glossaries}
\usepackage{glossary-superragged}
\setglossarystyle{superraggedheaderborder}
\usepackage{setspace}
\usepackage{parskip}

% package content table
\usepackage{tocbasic}

\usepackage{blindtext}

% ===================================================

% \renewcommand{\bibname}{Danh_sach_tai_lieu_tham_khao} 
\usepackage[backend=bibtex,style=ieee]{biblatex}  %backend=biber is 'better'

\addbibresource{Chapter/reference.bib} % chèn file chứa danh mục tài liệu tham khảo vào 

\include{lstlisting} % Phần này cho phép chèn code và formatting code như C, C++, Python

\setglossarystyle{longheaderborder} % Style có khung và tiêu đề
% Tùy chỉnh lại độ rộng cột để không bị tràn lề
\renewenvironment{theglossary}%
  {\begin{longtable}{|p{0.2\textwidth}|p{0.7\textwidth}|}}% p{0.7\textwidth} giúp tự xuống dòng
  {\end{longtable}}
\makenoidxglossaries
% List of Abbreviations and Acronyms
\newglossaryentry{OOP}{
    type=\acronymtype,
    name={OOP},
    description={Object-Oriented Programming},
    first={Object-Oriented Programming (OOP)}
}
\newglossaryentry{GUI}{
    type=\acronymtype,
    name={GUI},
    description={Graphical User Interface},
    first={Graphical User Interface (GUI)}
}
\newglossaryentry{ITS}{
    type=\acronymtype,
    name={ITS},
    description={Intelligent Transportation Systems},
    first={Intelligent Transportation Systems (ITS)}
}


% ===================================================
%Thay đổi font của công thước toán học - dùng Xe hoặc Lua
% \usepackage{unicode-math}
% \setmathfont{CambriaMath} %[Extension={.ttf},Path=./CambriaFonts/,Scale=1]
% \setmathrm{Cambria} %[Extension={.ttf},Path=./CambriaFonts/,Scale=1]
% \usepackage{fontspec}
% \setmainfont{Times New Roman}
\sloppy



\fancypagestyle{plain}{%
\fancyhf{} % clear all header and footer fields
\fancyfoot[RO,RE]{\thepage} %RO=right odd, RE=right even
\renewcommand{\headrulewidth}{0pt}
\renewcommand{\footrulewidth}{0pt}}

\setlength{\headheight}{10pt}

\def \TITLE{GRADUATION THESIS}
\def \AUTHOR{Phùng Minh Phúc}

% ===================================================
\titleformat{\chapter}[hang]{\centering\bfseries}{CHAPTER \thechapter.\ }{0pt}{}[]

\titleformat 
    {\chapter} % command
    [hang] % shape
    {\centering\bfseries} % format
    {CHAPTER \thechapter.\ } % label
    {0pt} %sep
    {} % before
    [] % after
\titlespacing*{\chapter}{0pt}{-20pt}{20pt}

\titleformat
    {\section} % command
    [hang] % shape
    {\bfseries} % format
    {\thechapter.\arabic{section}\ \ \ \ } % label
    {0pt} %sep
    {} % before
    [] % after
\titlespacing{\section}{0pt}{\parskip}{0.5\parskip}

\titleformat
    {\subsection} % command
    [hang] % shape
    {\bfseries} % format
    {\thechapter.\arabic{section}.\arabic{subsection}\ \ \ \ } % label
    {0pt} %sep
    {} % before
    [] % after
\titlespacing{\subsection}{30pt}{\parskip}{0.5\parskip}

\renewcommand\thesubsubsection{\alph{subsubsection}}
\titleformat
    {\subsubsection} % command
    [hang] % shape
    {\bfseries} % format
    {\alph{subsubsection}, \ } % label
    {0pt} %sep
    {} % before
    [] % after
\titlespacing{\subsubsection}{50pt}{\parskip}{0.5\parskip}

% \newcommand{\titlesize}{\fontsize{18pt}{23pt}\selectfont}
% \newcommand{\subtitlesize}{\fontsize{16pt}{21pt}\selectfont}
% \titleclass{\part}{top}
% \titleformat{\part}[display]
%   {\normalfont\huge\bfseries}{\centering}{20pt}{\Huge\centering}
% \titlespacing{\part}{0pt}{em}{1em}
% \titlespacing{\section}{0pt}{\parskip}{0.5\parskip}
% \titlespacing{\subsection}{0pt}{\parskip}{0.5\parskip}
% \titlespacing{\subsubsection}{0pt}{\parskip}{0.5\parskip}



% ===================================================
\usepackage{hyperref}
\hypersetup{pdfborder = {0 0 0}}
\hypersetup{pdftitle={\TITLE},
	pdfauthor={\AUTHOR}}
	
\usepackage[all]{hypcap} % Cho phép tham chiếu chính xác đến hình ảnh và bảng biểu

\graphicspath{{figures/}{../figures/}} % Thư mục chứa các hình ảnh

\counterwithin{figure}{chapter} % Đánh số hình ảnh kèm theo chapter. Ví dụ: Hình 1.1, 1.2,..

\title{\bf \TITLE}
\author{\AUTHOR}

\setcounter{secnumdepth}{3} % Cho phép subsubsection trong report
% \setcounter{tocdepth}{3} % Chèn subsubsection vào bảng mục lục

\theoremstyle{definition}
\newtheorem{example}{Example}[chapter] % Định nghĩa môi trường ví dụ

\onehalfspacing
\setlength{\parskip}{6pt}
\setlength{\parindent}{15pt}



% =========================== BODY ===============
\begin{document}
% \newgeometry{top=2cm, bottom=2cm, left=2cm, right=2cm}
\subfile{Chapter/Cover} % Phần bìa
% \subfile{Chapter/Cover2} % Phần bìa lót
% \restoregeometry

% ===================================================
\pagestyle{empty} % Header và footer rỗng
%\newpage
%\pagenumbering{gobble} % Xóa page numbering ở cuối trang
%\subfile{chapters/0_1_subject.tex}

% \pagestyle{empty} % Header và footer rỗng
\newpage
\pagenumbering{gobble} % Xóa page numbering ở cuối trang
\subfile{Chapter/0_2_Acknowledgment.tex}

% \pagestyle{empty} % Header và footer rỗng
\newpage
\pagenumbering{gobble} % Xóa page numbering ở cuối trang
\subfile{Chapter/0_3_Abstract.tex}


% ===================================================
% \pagestyle{empty} % Header và footer rỗng
\newpage
\pagenumbering{gobble} % Xóa page numbering ở cuối trang
\renewcommand*\contentsname{TABLE OF CONTENTS}
\titlecontents{chapter}
    [0.0cm]             % left margin
    {\bfseries\vspace{0.3cm}}                  % above code
    {{\bfseries{\scshape} CHAPTER \thecontentslabel.\ }} % numbered format
    {}         % unnumbered format
    {\titlerule*[0.3pc]{.}\contentspage}         % filler-page-format, e.g dots
    
\titlecontents{section}
    [0.0cm]             % left margin
    {\vspace{0.3cm}}                  % above code
    {\thecontentslabel \ } % numbered format
    {}         % unnumbered format
    {\titlerule*[0.3pc]{.}\contentspage}         % filler-page-format, e.g dots
    
\titlecontents{subsection}
    [1.0cm]             % left margin
    {\vspace{0.3cm}}                  % above code
    {\thecontentslabel \ } % numbered format
    {}         % unnumbered format
    {\titlerule*[0.3pc]{.}\contentspage}         % filler-page-format, e.g dots

 % Tạo mục lục tự động
\addtocontents{toc}{\protect\thispagestyle{empty}}
\tableofcontents 
\thispagestyle{empty}
\cleardoublepage

\pagenumbering{roman}
%Tạo danh mục hình vẽ.
\renewcommand{\listfigurename}{LIST OF FIGURES}
{\let\oldnumberline\numberline
\renewcommand{\numberline}{Figure~\oldnumberline}
\listoffigures} 
% \phantomsection\addcontentsline{toc}{section}{\numberline {} DANH MỤC HÌNH VẼ}
\newpage


 %Tạo danh mục bảng biểu.
\renewcommand{\listtablename}{LIST OF TABLES}
{\let\oldnumberline\numberline
\renewcommand{\numberline}{Table~\oldnumberline}
\listoftables}
% \phantomsection\addcontentsline{toc}{section}{\numberline {} DANH MỤC BẢNG BIỂU}

\glsaddall 
% \renewcommand*{\glossaryname}{Danh sách thuật ngữ}
\renewcommand*{\acronymname}{LIST OF ABBREVIATIONS}
\renewcommand*{\entryname}{Abbreviation}
\renewcommand*{\descriptionname}{Definition}

% In danh sách
\printnoidxglossary[type=\acronymtype] % Chỉ định in loại từ viết tắt (acronym)
% \phantomsection\addcontentsline{toc}{section}{\numberline {} DANH MỤC THUẬT NGỮ VÀ TỪ VIẾT TẮT}

% \renewcommand\appendixname{APPENDIX}
% \renewcommand\appendixpagename{APPENDIX}
% \renewcommand\appendixtocname{APPENDIX}

\renewcommand{\figurename}{Figure}
\renewcommand{\tablename}{Table}
\renewcommand{\chaptername}{CHAPTER}

% ===================================================


\newpage
\pagenumbering{arabic}

\pagestyle{fancy}
\fancyhf{}
\fancyhead[RE, LO]{\leftmark}
%\fancyhead[LE]{\rightmark}
\fancyfoot[RE, LO]{\thepage}

\chapter{INTRODUCTION}
\subfile{Chapter/1_Introduction} % Phần mở đầu

\newpage
%\pagestyle{fancy} % Áp dụng header và footer
\chapter{METHODOLOGY}
\subfile{Chapter/2_Methodology}

\newpage
% \pagestyle{fancy} % Áp dụng header và footer
\chapter{RESULTS DISCUSSION}
\subfile{Chapter/3_Numerical_results}

\newpage
\chapter{CONCLUSIONS}
\subfile{Chapter/4_Conclusions}

\newpage
% \pagestyle{fancy} % Áp dụng header và footer
% \chapter*{SHORT NOTICES ON REFERENCE} %Kết luận và hướng phát triển}
% \label{chapter:reference}
% \subfile{Chapter/5_Reference}

% ===================================================
\newpage
\renewcommand\bibname{REFERENCE}
\printbibliography
\phantomsection\addcontentsline{toc}{chapter}{REFERENCE}

% \appendixpage
% \appendix
% \addappheadtotoc

\titleformat{\chapter}[hang]{\centering\bfseries}{ \thechapter.\ }{0pt}{}[]
\titlespacing*{\chapter}{0pt}{-20pt}{20pt}

\titlecontents{chapter}
    [0.0cm]             % left margin
    {\bfseries\vspace{0.3cm}}                  % above code
    {{\bfseries{\scshape} \thecontentslabel.\ }} % numbered format
    {}         % unnumbered format
    {\titlerule*[0.3pc]{.}\contentspage} 
    
% \chapter{GRADUATION THESIS GUIDANCE}
% \subfile{Chapter/Appendix_A}

% \newpage
% \chapter{USE CASE DESCRIPTIONS}
% \subfile{Chapter/Appendix_B}

\end{document}
